\documentclass[10pt,a4paper]{amsart}
\usepackage{amsmath,amssymb,amsthm,graphicx}
% Retained preamble for the Stacks Project chapter morphisms.tex (Morphisms of Schemes).
% The chapter file pinned-original-source.tex carries only \input{preamble} on its first line, so no
% preamble text lives inside the chapter itself.  The Stacks Project uses one single shared preamble for
% every chapter; the snapshot stages it verbatim in sources/stacks-chow/preamble.tex.  The block below is
% copied mechanically, line for line, from that file's lines 14-18 (the Xy-pic setup) and lines 154-217
% (the native theorem declarations together with the mathematical macros).  Nothing was added, removed or
% redefined.  The wrapper class and the amsmath/amssymb/amsthm/graphicx base packages are supplied by the
% assembler, so the class/\externaldocument part of the shared preamble is not repeated here.
% For commutative diagrams we use Xy-pic
\usepackage[all]{xy}

% We use 2cell for 2-commutative diagrams.
\xyoption{2cell}
% Theorem environments.
%
\theoremstyle{plain}
\newtheorem{theorem}[subsection]{Theorem}
\newtheorem{proposition}[subsection]{Proposition}
\newtheorem{lemma}[subsection]{Lemma}

\theoremstyle{definition}
\newtheorem{definition}[subsection]{Definition}
\newtheorem{example}[subsection]{Example}
\newtheorem{exercise}[subsection]{Exercise}
\newtheorem{situation}[subsection]{Situation}

\theoremstyle{remark}
\newtheorem{remark}[subsection]{Remark}
\newtheorem{remarks}[subsection]{Remarks}

\numberwithin{equation}{subsection}

% Macros
%
\def\lim{\mathop{\mathrm{lim}}\nolimits}
\def\colim{\mathop{\mathrm{colim}}\nolimits}
\def\Spec{\mathop{\mathrm{Spec}}}
\def\Hom{\mathop{\mathrm{Hom}}\nolimits}
\def\Ext{\mathop{\mathrm{Ext}}\nolimits}
\def\SheafHom{\mathop{\mathcal{H}\!\mathit{om}}\nolimits}
\def\SheafExt{\mathop{\mathcal{E}\!\mathit{xt}}\nolimits}
\def\Sch{\mathit{Sch}}
\def\Mor{\mathop{\mathrm{Mor}}\nolimits}
\def\Ob{\mathop{\mathrm{Ob}}\nolimits}
\def\Sh{\mathop{\mathit{Sh}}\nolimits}
\def\NL{\mathop{N\!L}\nolimits}
\def\CH{\mathop{\mathrm{CH}}\nolimits}
\def\proetale{{pro\text{-}\acute{e}tale}}
\def\etale{{\acute{e}tale}}
\def\QCoh{\mathit{QCoh}}
\def\Ker{\mathop{\mathrm{Ker}}}
\def\Im{\mathop{\mathrm{Im}}}
\def\Coker{\mathop{\mathrm{Coker}}}
\def\Coim{\mathop{\mathrm{Coim}}}

% Boxtimes
%
\DeclareMathSymbol{\boxtimes}{\mathbin}{AMSa}{"02}

%
% Macros for moduli stacks/spaces
%
\def\QCohstack{\mathcal{QC}\!\mathit{oh}}
\def\Cohstack{\mathcal{C}\!\mathit{oh}}
\def\Spacesstack{\mathcal{S}\!\mathit{paces}}
\def\Quotfunctor{\mathrm{Quot}}
\def\Hilbfunctor{\mathrm{Hilb}}
\def\Curvesstack{\mathcal{C}\!\mathit{urves}}
\def\Polarizedstack{\mathcal{P}\!\mathit{olarized}}
\def\Complexesstack{\mathcal{C}\!\mathit{omplexes}}
% \Pic is the operator that assigns to X its picard group, usage \Pic(X)
% \Picardstack_{X/B} denotes the Picard stack of X over B
% \Picardfunctor_{X/B} denotes the Picard functor of X over B
\def\Pic{\mathop{\mathrm{Pic}}\nolimits}
\def\Picardstack{\mathcal{P}\!\mathit{ic}}
\def\Picardfunctor{\mathrm{Pic}}
\def\Deformationcategory{\mathcal{D}\!\mathit{ef}}

% --- appended 2026-10-12 (Chief): the project own definitions of the reference/slogan/history
% display environments, copied mechanically from the shared Stacks preamble lines 8-12
% (\usepackage{verbatim} plus the three \newenvironment lines, which the chapter preamble keeps
% verbatim from the project preamble). Nothing added, removed or redefined: these environments
% carry no mathematical content and the project itself typesets them as comments. Without these
% three lines a counted statement whose own text uses them fails with Environment slogan undefined.
\usepackage{verbatim}
\newenvironment{reference}{\comment}{\endcomment}
\newenvironment{slogan}{\comment}{\endcomment}
\newenvironment{history}{\comment}{\endcomment}
\usepackage{cleveref}
\crefname{originallemma}{original lemma}{original lemmas}
\Crefname{originallemma}{Original lemma}{Original lemmas}
\makeatletter
\newlabel{algebra-lemma-characterize-finite-presentation}{{00QO}{source}{}{sourcepointer-1}{}}
\newlabel{algebra-lemma-characterize-finite-presentation@cref}{{[originallemma][0][]00QO}{[1][0][]source}{}{}{}}
\makeatother
\crefname{originallemma}{original lemma}{original lemmas}
\Crefname{originallemma}{Original lemma}{Original lemmas}
\makeatletter
\newlabel{algebra-lemma-localization-colimit}{{00CR}{source}{}{sourcepointer-2}{}}
\newlabel{algebra-lemma-localization-colimit@cref}{{[originallemma][0][]00CR}{[1][0][]source}{}{}{}}
\makeatother
\crefname{originallemma}{original lemma}{original lemmas}
\Crefname{originallemma}{Original lemma}{Original lemmas}
\makeatletter
\newlabel{algebra-lemma-subring-of-local-ring}{{0BX1}{source}{}{sourcepointer-3}{}}
\newlabel{algebra-lemma-subring-of-local-ring@cref}{{[originallemma][0][]0BX1}{[1][0][]source}{}{}{}}
\makeatother
\title{A morphism is determined on a neighbourhood by its map on local rings, and a local map of local rings is realized by a morphism, in four standard situations}
\author{The Stacks Project Authors}
\date{stacks-project:morphisms.tex:local-snapshot-20261012; GNU FDL 1.2 or later}
\begin{document}
\maketitle
\noindent\textit{Editorial note.} All delivered mathematical text is the author's own, copied from the pinned Stacks Project chapter file pinned-original-source.tex; no statement or proof step was rewritten, reordered or removed.\par\medskip
\noindent\textit{Editorial note.} Following the multi-claim convention for this source, source\_url is the source-level URL https://stacks.math.columbia.edu/ and the source label lemma-morphism-defined-local-ring is carried in claim\_key and here rather than in a per-claim tag URL.\par\medskip
\noindent\textit{Editorial note.} The preamble is the chapter's shared Stacks Project preamble, retained as sources/stacks-morphisms/retained-preamble.tex: a line-for-line excerpt of the shared Stacks preamble staged with this snapshot. No macro was invented or redefined, and this card adds no preamble text of its own.\par\medskip
\noindent\textit{Editorial note.} Statement span 10338-10366 and proof span 10368-10422 come from the halfA pool card operations/pool-cards/stacks-morphisms-halfA.json and were re-read in the pinned chapter file before this card was written.\par\medskip
\noindent\textit{Editorial note.} The retained author text contains a typographical slip in its last sentence ('by another application of See Algebra, Lemma ...'). Under the lane rule that all delivered mathematical text is copied from the pinned chapter file without rewriting, the sentence is retained verbatim and no correction is made.\par\medskip
\noindent\textit{Editorial note.} Tag self-check per LANE\_BRIEF 10.5 (5): the label lemma-morphism-defined-local-ring resolves in the project's own tag list to canonical tag 0BX6 (line 14969 of the index, full label morphisms-lemma-morphism-defined-local-ring). reference/public-items.jsonl holds no row whose primary source\_file is morphisms.tex, so the intersection with the same-chapter public tag ids is empty and the candidate passes.\par\medskip
\noindent\textit{Editorial note.} The gate was run as: python operations/agent-authoring/tools/precard\_gate.py --source morphisms.tex --label lemma-morphism-defined-local-ring --start 10338 --end 10366 --proof-start 10368 --proof-end 10422 --local-id 34091, giving RESULT=OK with all four span guards passing, no public interval row for this chapter and occupancy FREE for morphisms.tex 10338 10422. Spans were recomputed live from the label in the pinned chapter file; the located text matched the pool byte counts (statement 1038 characters, proof 2152 characters after whitespace removal). The proof span contains a commutative diagram drawn with xymatrix, which the retained 88-line chapter preamble supports.\par\medskip
\noindent\textit{Editorial note.} The candidate was newly offered by the queue only after the handed-over windows were filtered out of the queue helper, so the counting region of this card is outside the windows 8000-10000 and 12000-14000 of this chapter that were transferred to lane-s1.\par\medskip
\noindent\textit{Editorial note.} 3 cross-references occur inside the counted spans; they are listed in external\_references as navigation pointers only.\par\medskip
\noindent\textit{Allowed standard prerequisites (editorial).}\begin{itemize}
\item A subring of a local ring contains the localization of the subring at an element outside the maximal ideal, cited from the algebra chapter. Cited at proof lines 10388-10389 and 10419-10420. Produces the element g outside the relevant prime over which the two maps agree, and later the element over which a map from the base into the local ring factors.
\item The characterization of finite presentation, and the description of a localization as a filtered colimit, cited from the algebra chapter. Cited at proof lines 10410-10412. Lets the given local map from a finitely presented algebra be realized, after localizing the source at one element, by an honest algebra map, which is what produces the required morphism in the first case of part (2).
\end{itemize}\medskip
\ifdefined\hypertarget\hypertarget{sourcepointer-1}{}\fi
\noindent\textit{Source reference.} Original lemma 00QO: focused-seeds/AGP/sources/stacks-snapshot-20261011/tags.txt, source lines 974--974. Same-chapter reference used in the retained author text at source line(s) 10408, 11936. It is supplied only as a navigation pointer so that the retained text compiles; the counted item delivers the author's own statement and the author's complete proof, and no part of the referenced item is copied in or substituted for the counted argument.\par\medskip
\ifdefined\hypertarget\hypertarget{sourcepointer-2}{}\fi
\noindent\textit{Source reference.} Original lemma 00CR: focused-seeds/AGP/sources/stacks-snapshot-20261011/tags.txt, source lines 471--471. Same-chapter reference used in the retained author text at source line(s) 10409. It is supplied only as a navigation pointer so that the retained text compiles; the counted item delivers the author's own statement and the author's complete proof, and no part of the referenced item is copied in or substituted for the counted argument.\par\medskip
\ifdefined\hypertarget\hypertarget{sourcepointer-3}{}\fi
\noindent\textit{Source reference.} Original lemma 0BX1: focused-seeds/AGP/sources/stacks-snapshot-20261011/tags.txt, source lines 14964--14964. Same-chapter reference used in the retained author text at source line(s) 10389, 10421. It is supplied only as a navigation pointer so that the retained text compiles; the counted item delivers the author's own statement and the author's complete proof, and no part of the referenced item is copied in or substituted for the counted argument.\par\medskip
\begin{lemma}
\label{lemma-morphism-defined-local-ring}
Let $S$ be a scheme. Let $X$, $Y$ be schemes over $S$.
Let $s \in S$ and $x \in X$, $y \in Y$ points over $s$.
\begin{enumerate}
\item Let $f, g : X \to Y$ be morphisms over $S$ such that
$f(x) = g(x) = y$ and
$f^\sharp_x = g^\sharp_x : \mathcal{O}_{Y, y} \to \mathcal{O}_{X, x}$.
Then there is an open neighbourhood $U \subset X$ with
$f|_U = g|_U$ in the following cases
\begin{enumerate}
\item $Y$ is locally of finite type over $S$,
\item $X$ is integral,
\item $X$ is locally Noetherian, or
\item $X$ is reduced with finitely many irreducible components.
\end{enumerate}
\item Let $\varphi : \mathcal{O}_{Y, y} \to \mathcal{O}_{X, x}$
be a local $\mathcal{O}_{S, s}$-algebra map. Then there exists
an open neighbourhood $U \subset X$ of $x$ and a morphism $f : U \to Y$
mapping $x$ to $y$ with $f^\sharp_x = \varphi$ in the following cases
\begin{enumerate}
\item $Y$ is locally of finite presentation over $S$,
\item $Y$ is locally of finite type and $X$ is integral,
\item $Y$ is locally of finite type and $X$ is locally Noetherian, or
\item $Y$ is locally of finite type and $X$ is reduced with finitely
many irreducible components.
\end{enumerate}
\end{enumerate}
\end{lemma}

\begin{proof}
Proof of (1). We may replace $X$, $Y$, $S$ by suitable affine open
neighbourhoods of $x$, $y$, $s$ and reduce to the following algebra
problem: given a ring $R$, two $R$-algebra maps $\varphi, \psi : B \to A$
such that
\begin{enumerate}
\item $R \to B$ is of finite type, or $A$ is a domain, or $A$
is Noetherian, or $A$ is reduced and has finitely many minimal primes,
\item the two maps $B \to A_\mathfrak p$ are the same for
some prime $\mathfrak p \subset A$,
\end{enumerate}
show that $\varphi, \psi$ define the same map $B \to A_g$ for
a suitable $g \in A$, $g \not \in \mathfrak p$. If $R \to B$ is of
finite type, let $t_1, \ldots, t_m \in B$ be generators of $B$
as an $R$-algebra. For each $j$ we can find
$g_j \in A$, $g_j \not \in \mathfrak p$
such that $\varphi(t_j)$ and $\psi(t_j)$ have the same image in
$A_{g_j}$. Then we set $g = \prod g_j$.
In the other cases (if $A$ is a domain, Noetherian, or reduced
with finitely many minimal primes), we can find a $g \in A$,
$g \not \in \mathfrak p$ such that $A_g \subset A_\mathfrak p$.
See Algebra, Lemma \ref{algebra-lemma-subring-of-local-ring}.
Thus the maps $B \to A_g$ are equal as desired.

\medskip\noindent
Proof of (2). To do this we may replace $X$, $Y$, and $S$ by suitable affine
opens. Say $X = \Spec(A)$, $Y = \Spec(B)$, and $S = \Spec(R)$.
Let $\mathfrak p \subset A$ be the prime ideal corresponding to $x$.
Let $\mathfrak q \subset B$ be the prime corresponding to $y$.
Then $\varphi$ is a local $R$-algebra map
$\varphi : B_\mathfrak q \to A_\mathfrak p$.
If $R \to B$ is a ring map of finite presentation, then there exists a
$g \in A \setminus \mathfrak p$ and an $R$-algebra map $B \to A_g$ such that
$$
\xymatrix{
B_\mathfrak q \ar[r]_\varphi & A_\mathfrak p \\
B \ar[u] \ar[r] & A_g \ar[u]
}
$$
commutes, see
Algebra, Lemmas \ref{algebra-lemma-characterize-finite-presentation} and
\ref{algebra-lemma-localization-colimit}.
The induced morphism $\Spec(A_g) \to \Spec(B)$ works.
If $B$ is of finite type over $R$, let $t_1, \ldots, t_m \in B$ be
generators of $B$ as an $R$-algebra. Then we can choose
$g_j \in A$, $g_j \not \in \mathfrak p$
such that $\varphi(t_j) \in \Im(A_{g_j} \to A_\mathfrak p)$.
Thus after replacing $A$ by $A[1/\prod g_j]$ we may assume
that $B$ maps into the image of $A \to A_\mathfrak p$.
If we can find a $g \in A$, $g \not \in \mathfrak p$
such that $A_g \to A_\mathfrak p$ is injective, then
we'll get the desired $R$-algebra map $B \to A_g$.
Thus the proof is finished by another application of
See Algebra, Lemma \ref{algebra-lemma-subring-of-local-ring}.
\end{proof}

\end{document}
